\section{Attained energies from restricted residuals}
\label{sec:ground_energy_residual}

A normalized eigenvector of the perturbed operator belongs to its actual
low-energy span. The restricted residual therefore bounds its energy
below a lower energy of the original operator, without a gap assumption.

\begin{lemma}[An eigenvector below the cutoff]
    \label{lem:eigenvector_mem_low_span}
    \lean{Submodule.mem_spectralSpanLE_of_eigenvector}
    \leanok
    \uses{def:spectral_span_le}
    Let $b_j$ be a finite complete orthonormal eigenbasis of a linear
    operator $B$, with real eigenvalues $a_j$. If $Bx=\mu x$ and
    $\mu\leq c$, then $x\in T_B(c)$.
\end{lemma}
\begin{proof}\leanok
    Taking the $j$th coefficient in $Bx=\mu x$ gives
    $(\mu-a_j)\langle b_j,x\rangle=0$.
    Every coefficient with $a_j>c$ therefore vanishes.
\end{proof}

\begin{lemma}[A lower energy and an attained eigenvalue]
    \label{lem:ground_energy_restricted_residual}
    \lean{real_eigenvalue_ge_of_unit_residual,
        ContinuousLinearMap.energy_sub_eigenvalue_le_norm_domRestrict,
        Matrix.energy_sub_eigenvalue_le_norm_spectral_residual}
    \leanok
    \uses{lem:eigenvector_mem_low_span}
    Let $A,B$ be maps on a complex inner-product space, $\|x\|=1$,
    $Bx=\mu x$ and
    $e\|x\|^2\leq\operatorname{Re}\langle x,Ax\rangle$.
    If $\|Ax-Bx\|\leq\xi$, then $e-\xi\leq\mu$.

    Suppose moreover that $A,B$ are continuous linear maps and $B$
    has a finite complete orthonormal eigenbasis with real eigenvalues.
    For every $c\geq\mu$,
    \begin{align}
        e-\mu\leq\|(A-B)|_{T_B(c)}\|.
    \end{align}
    For finite complex matrices, the same assertion follows from
    $A-eI\geq0$ in place of the quadratic-form lower bound.
\end{lemma}
\begin{proof}\leanok
    Since $\operatorname{Re}\langle x,Bx\rangle=\mu$,
    \begin{align}
        e-\mu
        \leq\operatorname{Re}\langle x,Ax-Bx\rangle
        \leq\|x\|\,\|Ax-Bx\|\leq\xi.
    \end{align}
    For the restricted assertion,
    Lemma~\ref{lem:eigenvector_mem_low_span} gives $x\in T_B(c)$.
    The operator norm of the restriction bounds its value at this
    unit vector. Finally, positivity of $A-eI$ gives the required
    quadratic-form inequality in the matrix case.
\end{proof}

These are auxiliary estimates for the soft-truncation spectral
comparison in
\cite[lines 1195--1255 and 1373--1392]{Arad2016RigorousRG}.
No residual bound or spectral gap is inferred without the corresponding
hypotheses or physical construction.
